\section{Forschungs- und Entwicklungsperspektive}
\label{sec:08-03_future-research}

Die in dieser Arbeit gewonnenen Erkenntnisse liefern belastbare Hinweise auf die Wirkmechanismen nutzerzentrierter \Gls{ia} und \Gls{uxstrategien} im organisationalen Kontext. Zugleich zeigen die identifizierten Limitationen konkrete Ansatzpunkte für weiterführende Forschung und die praxisorientierte Weiterentwicklung vergleichbarer Wissensplattformen.

Eine naheliegende Forschungsperspektive besteht in der Ausweitung des Untersuchungsdesigns auf mehrere Organisationen und Plattformtypen. Vergleichende Mehrfallstudien könnten prüfen, inwieweit die beobachteten Effekte von struktureller Klarheit, Einstiegslogiken und Governance auch in unterschiedlichen Branchen, Reifegraden und Systemlandschaften stabil auftreten. Dadurch ließen sich kontextspezifische Einflussfaktoren differenzierter erfassen und die externe Validität der Befunde erhöhen. Ergänzend bieten sich quantitative Studien mit größeren Stichproben an, um belastbare Aussagen zu Effektstärken und Zusammenhängen zwischen Strukturqualität, Nutzungseffizienz und Akzeptanz zu ermöglichen.

Ein weiterer Forschungsbedarf ergibt sich im Bereich longitudinaler Analysen. Die vorliegende Arbeit erfasst primär kurzfristige Nutzungseffekte und subjektive Einschätzungen nach der Einführung des neuen Wikis. Zukünftige Studien könnten untersuchen, wie sich Nutzungsmuster, Akzeptanz und wahrgenommener Nutzen über längere Zeiträume entwickeln und ob anfängliche Effizienzgewinne stabil bleiben oder sich weiter verstärken. Die Kombination von Nutzungslogdaten, qualitativen Folgeinterviews und wiederholten \Gls{usability}-Messungen würde eine differenzierte Betrachtung von Lernprozessen, Routinenbildung und organisationaler Verankerung ermöglichen.

Auch methodisch bieten sich Erweiterungen an. Eine vollständige Transkription der Interviews sowie die stärkere Einbindung mehrerer Kodierender könnten die Nachvollziehbarkeit qualitativer Analysen weiter erhöhen. Darüber hinaus könnten experimentellere Designs eingesetzt werden, um einzelne Gestaltungsvariablen -- etwa Navigationslogiken, Einstiegsseiten oder Suchmechanismen -- gezielter zu untersuchen und kausale Effekte präziser zu bestimmen.

Inhaltlich eröffnet sich weiteres Forschungspotenzial im Zusammenspiel von \Gls{ia}, Suchsystemen und algorithmischer Unterstützung. Mit der zunehmenden Verbreitung semantischer Suche, Empfehlungssystemen und KI-gestützten Assistenzfunktionen stellt sich die Frage, wie klassische Strukturprinzipien der \Gls{ia} sinnvoll mit adaptiven Zugangsmechanismen kombiniert werden können, ohne Transparenz, Erwartungskonformität und Kontrollierbarkeit zu verlieren. Hier besteht Bedarf an empirischer Forschung zu Nutzungsvertrauen, Erklärbarkeit und kognitiver Belastung in hybriden Informationssystemen.

\newpage

Neben wissenschaftlichen Anschlussfragen ergeben sich auch konkrete Entwicklungsperspektiven für die untersuchte Plattform. Die Weiterentwicklung der Personen- und Ansprechpartnerlogik, die Optimierung der Suchfunktion sowie der Ausbau von \Gls{governance}-Strukturen bieten kurzfristige Ansatzpunkte zur Steigerung von Nutzungsqualität und Akzeptanz. Perspektivisch können datenbasierte Monitoring-Ansätze etabliert werden, um Nutzungsmuster kontinuierlich zu analysieren und strukturelle Anpassungen evidenzbasiert zu priorisieren. 

Ebenso bietet die schrittweise Integration kollaborativer Funktionen oder intelligenter Assistenzsysteme Potenzial, sofern diese konsequent an klarer Struktur, guter Auffindbarkeit und transparenter Governance ausgerichtet bleiben.

Auch auf organisationaler Ebene besteht weiteres Entwicklungspotenzial. Die Ergebnisse legen nahe, \Gls{ux} nicht nur als Gestaltungsdisziplin, sondern als kontinuierliche organisationale Lernpraxis zu verstehen. Zukünftige Projekte können untersuchen, wie sich entsprechende Kompetenzen, Rollen und Entscheidungsprozesse nachhaltig etablieren lassen und welche Auswirkungen dies auf Innovationsfähigkeit, Wissensarbeit und digitale Reife von Organisationen hat.

Insgesamt zeigt die vorliegende Arbeit, dass nutzerzentrierte \Gls{ia} ein wirksamer Hebel zur Verbesserung von Orientierung, Effizienz und Akzeptanz in organisationalen Wissensplattformen ist, ihre nachhaltige Wirkung jedoch erst im Zusammenspiel mit Governance, Kommunikation und organisationaler Verankerung entsteht. Damit leistet die Arbeit sowohl einen empirisch fundierten Beitrag zum Verständnis der Wirkmechanismen von \Gls{ia} und \Gls{uxstrategien} als auch eine praxisrelevante Orientierung für die Weiterentwicklung vergleichbarer Systeme. Die aufgezeigten Forschungs- und Entwicklungsperspektiven markieren zugleich Ansatzpunkte, um diese Erkenntnisse in zukünftigen Studien und Implementierungen weiter zu vertiefen.

